% Buffon's needle. Left panel: the geometry. A needle of length l dropped
% on ruled lines a distance d apart (l <= d) crosses a line when the
% distance x from the needle centre to the nearest line is at most
% (l/2) sin(theta). Right panel: the running Monte Carlo estimate of pi
% from the crossing fraction 2l/(d p_cross), stylised as a curve settling
% toward pi with 1/sqrt(n) fluctuations. Both panels are schematic.
\begin{figure}[htbp]
\centering
\begin{tikzpicture}[font=\small]

% Left panel: geometry.
\begin{scope}[x=1cm, y=1cm]
  % ruled lines at y = 0, 1.5, 3.0
  \foreach \y in {0, 1.5, 3.0} {
    \draw[enggray!55, thick] (0, \y) -- (4.4, \y);
  }
  % needle crossing a line (red)
  \draw[very thick, engred] (1.0, 1.1) -- (2.1, 1.85);
  \fill[engred] (1.55, 1.475) circle (1.2pt);
  % needle not crossing (blue)
  \draw[very thick, engblue] (2.7, 1.95) -- (3.7, 2.55);
  \fill[engblue] (3.2, 2.25) circle (1.2pt);
  % angle and offset annotations for the crossing needle
  \draw[dashed, engblack] (1.55, 1.475) -- (1.55, 1.5);
  \node[anchor=north east, engred, font=\scriptsize] at (2.1, 1.85)
    {crosses};
  \node[anchor=north west, engblue, font=\scriptsize] at (3.7, 2.55)
    {misses};
  \node[anchor=east, enggray, font=\scriptsize] at (0.0, 0.75) {$d$};
  \draw[<->, enggray] (0.15, 0.0) -- (0.15, 1.5);
  \node[anchor=south, engblack, font=\scriptsize] at (2.2, 3.05)
    {ruled plane};
\end{scope}

% Right panel: convergence of the estimate.
\begin{scope}[xshift=6.2cm]
\begin{axis}[
  width=7.4cm, height=5.6cm,
  xlabel={drops $n$ (log scale)},
  ylabel={estimate of $\pi$},
  xmode=log, xmin=10, xmax=100000,
  ymin=2.9, ymax=3.4,
  grid=both, grid style={enggray!18},
  tick label style={font=\scriptsize},
  label style={font=\scriptsize},
]
  % target line pi
  \addplot[very thick, engblack, domain=10:100000, samples=2]
    {3.14159265};
  % a stylised running estimate: pi plus a damped oscillation ~ 1/sqrt(n)
  \addplot[thick, engblue, domain=10:100000, samples=120]
    {3.14159265 + 0.9*sin(deg(1.7*ln(x)))/sqrt(x)};
  \node[anchor=south east, engblack, font=\scriptsize]
    at (axis cs:90000, 3.14159265) {$\pi$};
\end{axis}
\end{scope}

\end{tikzpicture}
\caption[Buffon's needle and its Monte Carlo estimate of $\pi$]{Left: a
needle of length $\ell$ dropped on lines a distance $d\ge\ell$ apart
crosses a line (red) when the offset of its centre from the nearest line
falls below $(\ell/2)\sin\theta$, and misses otherwise (blue). Averaging
over a uniform angle and a uniform offset gives crossing probability
$2\ell/(\pi d)$, so the crossing fraction estimates $\pi$. Right: a
running estimate of $\pi$ from the crossing fraction settles toward the
true value (black line) with fluctuations that shrink as $1/\sqrt{n}$,
the same square-root-of-$n$ wall every Monte Carlo estimator meets. The
needle turns a geometric probability into a physical computation of a
transcendental constant. Source: schematic by the authors; the crossing
probability is derived in section 13.1.}
\label{fig:vol02ch13:buffon-needle}
\end{figure}
